\documentclass[10pt,a4paper]{amsart}
\usepackage[margin=19mm]{geometry}
\usepackage{amsmath,amssymb,mathtools}
\usepackage[scr=rsfs,cal=euler]{mathalfa}
\usepackage{xfrac}
\usepackage[unicode,hidelinks,bookmarks=false]{hyperref}
\newcommand{\ie}{i.e.\ }
\newcommand{\cf}{cf.\ }
\newcommand{\resp}{respectively\ }
\newcommand{\Subsection}{\subsection}
\providecommand{\nobreakdash}{\nobreak\hbox{-}}
\setlength{\emergencystretch}{2em}
\multlinegap0pt
\usepackage{bm}
\usepackage{bbm}
\usepackage{dsfont}
\usepackage{xspace}
\usepackage{cancel}
\usepackage{mathtools}
\usepackage{amsthm}
\makeatletter
\@ifundefined{definition}{\newtheorem{definition}{Definition}}{}
\@ifundefined{lemma}{\newtheorem{lemma}[definition]{Lemma}}{}
\@ifundefined{proposition}{\newtheorem{proposition}[definition]{Proposition}}{}
\makeatother
\DeclareMathOperator{\Adm}{Adm}
\DeclareMathOperator{\SC}{sc}
\DeclareMathOperator{\pr}{pr}
\DeclareMathOperator{\rk}{rk}
\newcommand{\bC}{\mathbb{C}}
\newcommand{\cP}{\mathcal{P}}
\title[Delta-matroids and toric degenerations in OG(n,2n+1)]{Delta-matroids and toric degenerations in OG(n,2n+1)}
\author{Chen Chen, Carl Lian}
\date{Source: arXiv:2507.16133v2 (CC BY 4.0)}
\begin{document}
\maketitle

\noindent\emph{Source and licence.} Delta-matroids and toric degenerations in OG(n,2n+1). Version arXiv:2507.16133v2 (CC BY 4.0), distributed under the Creative Commons Attribution 4.0 International licence, as recorded in the arXiv licence metadata for this exact version and shown on the abstract page https://arxiv.org/abs/2507.16133v2. Full rights notice: licenses/arxiv-2507.16133-CC-BY-4.0.txt.\par\smallskip

\noindent\textit{Editorial scope.} Counted unit: the Proposition whose source label is \texttt{prop:S\_ineq\_positive} in the article (source0.tex lines 1557--1559, source label \texttt{prop:S\_ineq\_positive}), delivered together with the article's own complete original proof of it (source0.tex lines 1561--1620, one \texttt{proof} environment of 60 lines). No mathematics is written, repaired or re-derived: statement and proof are the article's own text line for line. Required context retained uncounted: prop:solve\_for\_* (lines 702--704); prop:solve\_for\_*\_saturated (lines 776--778); def:I\_associated\_to\_x (lines 1394--1406); prop:x\_in\_P(I) (lines 1475--1481); lem:split\_rank (lines 1503--1509). Every definition, display and label of those lines is retained unchanged. Omitted from this package and not redistributed: 1--701, 705--775, 779--1393, 1407--1474, 1482--1502, 1510--1556, 1560--1560, 1621--1967. Nothing inside a retained excerpt is omitted or altered: every retained body is delivered exactly as the article's lines are written, so no omission receipt of any kind accompanies this package. Citations: the retained excerpts carry no citation command at all, so no bibliography entry of the article is transcribed and no reference is redistributed. Reference adaptation: none was needed and none was made; every reference inside the delivered unit resolves inside it, so no unresolved reference or citation is produced. Printed slips kept literal: no printed word, symbol, hypothesis or number of the counted statement or of its proof was altered. Layout and class: the article's own class is \texttt{amsart} with packages including geometry, tikz, tikz-cd, url, hyperref, xy, amsmath, amssymb, enumitem, graphicx, mathdots, color, diagbox, array; the wrapper is the lane's pinned \texttt{amsart} editorial wrapper at 10pt on A4, which restates the theorem-like environments the retained text needs and adds the article's own macro definitions that the retained text needs (\texttt{\textbackslash Adm}, \texttt{\textbackslash SC}, \texttt{\textbackslash bC}, \texttt{\textbackslash cP}, \texttt{\textbackslash pr}, \texttt{\textbackslash rk}), with extra wrapper packages (bm, bbm, dsfont, xspace, cancel, mathtools). The delivered numbering is the unit's own. Assets: no retained span contains a figure environment, a graphics-inclusion command, a PDF-page inclusion, a table or a TikZ picture; no image file is copied into this package, the package declares no asset dependency and no image file is redistributed with it. Licence: version arXiv:2507.16133v2 of this article is distributed under the Creative Commons Attribution 4.0 International licence, as recorded in the arXiv licence metadata for this exact version and shown on the abstract page https://arxiv.org/abs/2507.16133v2. A full rights notice is delivered at licenses/arxiv-2507.16133-CC-BY-4.0.txt. Characters: every retained excerpt is pure ASCII, so the delivered engine, XeLaTeX with its default Latin Modern text font, reports no missing character.
\begin{proposition}\label{prop:solve_for_*}
The map $\pr$ is birational. That is, there is a non-empty Zariski open subset $V_{I,I'}\subset \bC^{d(I,I')}$ over which $\pr$ is an isomorphism.
\end{proposition}

\begin{proposition}\label{prop:solve_for_*_saturated}
    Proposition \ref{prop:solve_for_*} holds when $(I,I')$ is saturated. In fact, taking $V_{I,I'}\subset\bC^{d(I,I')}$ to be the locus where all of the $+$-entries are non-zero, $\pr$ is an isomorphism over $V_{I,I'}$. Furthermore, for any $A\in\pr^{-1}(V_{I,I'})$, all of the $*$-entries of $A$ are non-zero. 
\end{proposition}

\begin{definition}\label{def:I_associated_to_x}
    Let $x\in[0,1]^n$ be the point above. For all integers $\ell\in[1,\lceil y_n\rceil -1]$, let $\lambda_\ell\in[n-1]$ be the unique integer such that 
    \begin{equation}\label{eq:I'_inequalities}
       y_{n-\lambda_{\ell}}\le\ell< y_{n+1-\lambda_{\ell}}.
    \end{equation}
    Define the set
    \begin{equation*}
        \SC(x)=\{n+1-\lambda_1,\ldots,n+1-\lambda_{s}\}\subset\{2,\ldots,n\},
    \end{equation*}
    where $s=\max(0,\lceil y_n\rceil-1)$.
    
    Then, define $I=\{\lambda_1,\ldots,\lambda_{s}\}\subset[n-1]$. Define also $I'=\{\lambda'_1,\ldots,\lambda'_{s'}\}$, where $\lambda'_1>\cdots>\lambda_{s'}$, to be the complement of $I$ in $[n-1]$.
\end{definition}

\begin{proposition}\label{prop:x_in_P(I)}
    Let $x\in[0,1]^n$ be any point, and let $I\subset [n-1]$ be as in Definition \ref{def:I_associated_to_x}. Let $\Lambda_I\in\Sigma_{I,I^c}$ be a nice subspace. Then, for any $S\in\Adm_n$, we have
    \begin{equation*}
        x(S)\le g_{\Lambda_I}(S).
    \end{equation*}
    That is, we have $x\in\cP(\Lambda_I)$.
\end{proposition}

\begin{lemma}\label{lem:split_rank}
    Let $I \subset [n-1]$ and $\Lambda_I \in \Sigma_{I,I^c}$ be as in Proposition \ref{prop:x_in_P(I)}. Then,
    \begin{equation*}
        \rk_{\Lambda_I}(S) = \rk_{\Lambda_I}(S \cap [n])+ \rk_{\Lambda_I}(S \cap \overline{[n]})
    \end{equation*}
    for all $S \in \Adm_n$.
\end{lemma}

\begin{proposition}\label{prop:S_ineq_positive}
    Proposition \ref{prop:x_in_P(I)} holds when $S\subset[n]$.
\end{proposition}

\begin{proof}
    Let $\rho_1<\cdots<\rho_r$ be the indices such that
\begin{equation*}
    n +1 -\lambda_{\rho_1},n +1 -\lambda_{\rho_2}, \ldots,n +1 -\lambda_{\rho_r} \in \SC(x)
\end{equation*}
are the special columns of $M_{I,I^c}$ in $[n]$ \emph{not} contained in $S$. It will be convenient also to set
\begin{equation*}
    \rho_0=0,\lambda_0=n+1\text{ and }\rho_{r+1}=s+1,\lambda_{s+1}=0.
\end{equation*}

Then, for every integer $i \in \{0, \ldots, r\}$, define the set
\begin{equation*}
    \widetilde{S}_i =[(n+1-\lambda_{\rho_{i}})+1, (n+1-\lambda_{\rho_{i+1}})-1].
\end{equation*}
By definition, the $\widetilde{S}_i$ are disjoint. We also have 
\begin{equation*}
    \bigcup_{i \in \{0, \ldots, r\}} \widetilde{S}_i = [n] \setminus \{n +1 -\lambda_{\rho_1}, n +1 -\lambda_{\rho_2}, \ldots, n +1 -\lambda_{\rho_r}\}.
\end{equation*}
Note that the set $\widetilde{S}_i$ can be empty for any $i \in \{1, \ldots, r\}$. On the other hand, $1$ is always an element of $\widetilde{S}_0$. Taking $A$ to be as in the proof of Lemma \ref{lem:split_rank} above, we observe that the subsets $\widetilde{S}_i$ satisfy the property that 
\begin{equation*}
    \langle A_{\widetilde{S}_i}\rangle \cap \langle A_{\widetilde{S}_j}\rangle = 0
\end{equation*}
for any $i \neq j$. This implies that 
\begin{equation*}
    \rk_{\Lambda_I}\left(\bigcup_{i \in \{0, \ldots, r\}} \widetilde{S}_i\right) = \rk_{\Lambda_I}(\widetilde{S}_0) + \ldots + \rk_{\Lambda_I}(\widetilde{S}_r).
\end{equation*}
Now, we can write
\begin{equation*}
    S = S_0 \cup S_1 \cup \ldots \cup S_r,
\end{equation*}
where $S_i=S\cap \widetilde{S}_i$, so that
\begin{equation*}
    \rk_{\Lambda_I}(S) = \rk_{\Lambda_I}(S_0) + \ldots + \rk_{\Lambda_I}(S_r).
\end{equation*}
It now suffices to show that $x(S_i)\le g_{\Lambda_I}(S_i)= \rk_{\Lambda_I}(S_i)$ for each $i$. 

By construction, $S_i$ contains all of the $\rho_{i+1}-\rho_i-1$ special columns in the interval $\widetilde{S}_i=[(n+1-\lambda_{\rho_{i}})+1, (n+1-\lambda_{\rho_{i+1}})-1]$. Then, we have
\begin{equation*}
    g_{\Lambda_I}(S_i) = \rk_{\Lambda_I} (S_i) = 
    \begin{cases}
        \rho_{i+1}-\rho_i-1 &\text{ if } S_i = \widetilde{S}_i\cap \SC(x),\\
        \rho_{i+1}-\rho_i&\text{ otherwise}
    \end{cases}
\end{equation*}
Indeed, by the non-vanishing of the $*$-entries in Proposition \ref{prop:solve_for_*_saturated}, the column vectors in the set of special columns $\SC(x)\cap  \widetilde{S}_i$ are linearly independent, and if $S_i$ contains at least one additional column vector in a non-special column, then $\langle A_{S_i}\rangle=\langle A_{\widetilde{S}_i}\rangle$. In the first case, because $x\in[0,1]^n$, we have
\begin{equation*}
    x(S_i)\le |S_i|= \rho_{i+1}-\rho_i-1=\rk_{\Lambda_I} (S_i),
\end{equation*}
so the needed inequality is immediate.

Suppose, from now on, that there is at least one element in $S_i$ that is not a special column. Then, we have $\rk_{\Lambda_I}(S_i) = \rk_{\Lambda_I}(\widetilde{S}_i) = \rho_{i+1}-\rho_i$ and $x(S_i) \le x(\widetilde{S}_i)$. Thus, it suffices to show that 
\begin{equation*}
    x(\widetilde{S}_i) \le \rho_{i+1}-\rho_i,
\end{equation*}
for $i=0,\ldots,r$. On the other hand, we have
\begin{align*}
    x(\widetilde{S}_i)=y_{n-\lambda_{\rho_{i+1}}}-y_{n+1-\lambda_{\rho_i}}\le \rho_{i+1}-\rho_i,
\end{align*}
by Definition \ref{def:I_associated_to_x}. This completes the proof.
\end{proof}

\end{document}
